\section{Failure transparency}
\label{sec:failures}

The theorems of \S\ref{sec:theorems} cover failure-free runs. This
section adds server crashes, standby takeovers, Worker failures, retries,
zombie Workers, duplicate results, and StartFlow retries, and proves the
second theorem in the same fragment under the hypotheses of
\S\ref{sec:failure-scope}. For finite traces it gives two inclusions:
erasing the repeated external effects of a trace with failures gives a
failure-free trace for the word without retries, and every failure-free
trace is a trace with failures. For fair maximal traces with markers the
first inclusion holds for every admissible word, and the second for words
without StartFlow retries. Composed with the first theorem, it relates a
well-formed \Lang{} program to its execution on a failing server, where
each extra $\ev{out}(x)$ is known only to equal an earlier $\ev{out}(x)$
in value.

The failure model is Section~11 of the specification (revision~9), written
against the server at a later commit than the failure-free rules, again
with a citation for every rule. The failure-free rules are unchanged, so
the first proof stays valid. In Lean the model is \emph{DexF}
(\file{Failure/}); the failure-free model of the earlier sections is called
\emph{Dex} below.

\subsection{The failure model}
\label{sec:failure-model}

DexF is the Dex model plus four pieces of state and a set of failure rules.
The state records whether the server is up; a set $Z$ of \emph{zombie}
tasks, Worker tasks that the server gave up but that may still run; a set
$I$ of interrupted attempts, which recovery rebuilds; and, per invocation,
an attempt number, a phase Pending, Running, or Backoff, a task ID, and a
\emph{retained snapshot}. A Step method's attempts share one snapshot: the
server takes it at the first start and every retry reuses it.
Table~\ref{tab:failure-rules} lists the rules.

\begin{table}[t]
\centering
\small
\begin{tabularx}{\textwidth}{@{}P{4.1cm}L@{}}
\toprule
Rule & What happens \\
\midrule
\rulename{CRASH} & The server stops. Committed state, timers, the clock,
  and the zombies stay; every running task becomes a zombie; engine memory
  is lost. A log write with an unknown outcome is a crash. \\
\rulename{RECOVER} & The server restarts from the log. Every attempt that
  was Running joins $I$; it is never handed out again. \\
\rulename{RECOVER-INTERRUPT} & A separate commit per element of $I$: the
  attempt fails and goes to Backoff, or is ignored if a fence rejects it. \\
\rulename{ATTEMPT-LOST} & An attempt fails at any time (Worker error,
  method or heartbeat timeout, disconnect, server error). It keeps its
  snapshot, arms a retry timer, and its task becomes a zombie. \\
\rulename{ATTEMPT-EXHAUSTED} & The same failure when the retry policy
  allows no further attempt, or a non-retryable one: the Flow closes as
  Failed. \\
\rulename{ATTEMPT-RETURNED} & The Worker rejected the task; it goes back to
  Pending with the same attempt number and snapshot. \\
\rulename{ATTEMPT-HEARTBEAT} & A heartbeat checkpoint commits. \\
\rulename{RESULT-STALE} & A result, failure, or heartbeat whose fence (Step
  execution, method, attempt number, phase, task ID) no longer matches is
  ignored. Duplicate proposals are stale. \\
\rulename{RETRY-FIRE} & A due retry timer moves a Backoff invocation back
  to Pending, with its retained snapshot. It is part of
  \rulename{TIMER-FIRE}. \\
\rulename{ZOMBIE-EFFECT} & A zombie emits its next external effect, or is
  dropped. \\
\rulename{START-RETRY} & A StartFlow with the start's request ID and
  \kw{ignore\_already\_started\_error} returns the existing run. It is a
  stamped log entry and writes no run state. \\
\rulename{ENGINE-REJECT} & After the run was purged, the engine rejects a
  publication before proposing anything, with the event that apply's
  rejection would give. \\
\bottomrule
\end{tabularx}
\caption{The failure rules of DexF (specification \S11.3). They emit no
event, except \rulename{ATTEMPT-EXHAUSTED} ($\ev{closed}$),
\rulename{ZOMBIE-EFFECT} ($\ev{fx}$), \rulename{ENGINE-REJECT}
($\ev{rejected}$), and \rulename{RETRY-FIRE} when its command also fires
wait timers, which then acts as \rulename{TIMER-FIRE}.}
\label{tab:failure-rules}
\end{table}

Three choices make the model larger than the code, the safe direction
for an inclusion theorem: more DexF runs are more obligations. Failures
are nondeterministic and can happen at any time, with no model of how the
engine detects them. A zombie can emit effects at any later time, even
after a newer attempt of the same method has started. A zombie's result
can commit after a crash and before its interrupt, an order that the code
makes unlikely.

\paragraph{Durability.}
The durable point of a commit is its log batch in object storage, written
create-only. Recovery replays the same stamped commands through the same
apply function as live commits; it is not a special mode. A standby
takeover is a crash of the old owner followed by recovery on the new one.
The new owner writes an epoch marker at the first free log position, and
replay voids every later batch of the old epoch, so the log has one writer
at a time. The old owner may still carry frames to Workers until it stops;
those tasks are zombies, and their results cannot reach the new owner.

\subsection{Scope, decisions, and hypotheses}
\label{sec:failure-scope}

We fixed the scope before any proof (decisions D1 to D8 and OD1 of the
specification). Table~\ref{tab:failure-hypotheses} lists each assumption
and how the Lean model treats it.

\begin{table}[t]
\centering
\small
\begin{tabularx}{\textwidth}{@{}P{1.2cm}LP{4.6cm}@{}}
\toprule
 & Assumption & In Lean \\
\midrule
D1 & The fragment of the first theorem: one Flow, the admissible
  environment extended by StartFlow retries; no RPCs, SubFlows, StopFlow,
  or Attribute locks. & The Dex model, extended. \\
D2 & Every Step method is SYNC: an attempt reaches a Worker only after its
  start commits. An ASYNC start lost in a crash runs the same attempt again
  with a new snapshot. & \lean{IsSyncOnly}, true by construction: DexF has
  only the SYNC start. \\
H1 & No attempt exhausts its retries. Exhaustion closes the Flow as
  Failed, with WORKER\_API\_FAIL or INTERNAL in the code, a close that a
  failure-free run reaches only by deciding ForceFail. & Part of the trace
  sets; the model closes with the value \texttt{"WORKER\_API\_FAIL"}. \\
H2 & No failure policy (ProceedOnFailure, ProceedToConfiguredStep). &
  \lean{HasNoFailurePolicy}, true by construction. \\
H3 & User code ignores the attempt-varying context: attempt number, first
  attempt time, last error, last heartbeat value. Hard-delete Attributes,
  which are re-read on every attempt, are outside the fragment. & A
  hypothesis: \lean{IgnoresAttemptContext}. \\
H4 & No non-retryable failure. & A case of H1. \\
D6 & On one node, the wall clock does not step back between the timer
  service's due check and the stamping of its command. & The premise of
  the timer rule. \\
D7 & The input word is the \emph{outcome word}: requests in the order of
  their final outcomes. A re-sent publication is a second request. & The
  admissible outcome words. \\
HP & Every StartFlow retry is ordered in the log before the purge of the
  run. & True by construction: the model keeps a closed run. \\
TID & Every logged command of a started attempt (result, returned
  attempt, heartbeat, failure after start) carries a non-empty task ID. &
  True by construction: task IDs are minted numbers. \\
F0 & Every request of the outcome word gets its final outcome, in order.
  & \lean{RequestsArrive}, part of \lean{IsFairMaximalTrace}. \\
F1--F6 & The fairness of Table~\ref{tab:fairness}, restated for DexF:
  clocks grow without bound (F1); a timer of either kind that stays armed
  and due is fired (F2); an invocation that stays startable while the
  server is up starts (F3); every real task, in $\Theta$ or $Z$,
  eventually returns or is dropped (F4); every returned task eventually
  leaves by a Complete rule or a stale result (F5); F6 is trivial in the
  fragment. & \lean{DexF.IsFair}. \\
FF, FR, FI & For liveness: finitely many crashes, lost attempts, and
  returned attempts; the server comes back after every crash; every
  interrupt eventually leaves $I$, by its commit or by a crash that
  empties $I$, so after the last recovery each one commits. & Part of
  \lean{DexF.IsFair}. \\
OD1 & A timer command fires a non-empty prefix of the due timers in the
  order (fire time, encoded timer ID), as in the first theorem. & The
  timer rule of the Dex model for wait timers. Lean does not compute
  timer ID bytes: at equal fire times a wait timer comes before a retry
  timer, and retry timers keep their list order. \\
\bottomrule
\end{tabularx}
\caption{What the failure theorem assumes. H1 to H4 exclude four features
that change observable behavior beyond repeated effects.}
\label{tab:failure-hypotheses}
\end{table}

Four entries hold by construction. The Lean model has no failure policy,
no ASYNC start, and no purge of a closed run, and its task IDs are minted
numbers. The
theorem excludes those three features, and the named predicates make the
exclusion visible in every statement. H3 is the one real hypothesis on the
program. A DexF program is a function from the
attempt context to a Flow definition, so a program that reads its attempt
number is expressible, and the theorem excludes it explicitly.

D5, finitely many failures, is necessary: no system makes progress under
infinitely many crashes. D6 needs no clock
agreement between hosts. Stamps increase strictly after a takeover too,
because the new owner seeds its watermark from the largest applied stamp.
The timer order of OD1 is an idealization. The specification found one
window, present without any failure, in which the code fires a timer while
an earlier-due timer, armed by a command applied just before, is still
armed (Open question~26). We kept the prefix rule for both theorems so
that they compose; neither covers runs in that window.

\subsection{Repetitions and erasure}
\label{sec:repetitions}

DexF tags every $\ev{fx}$ event with a ghost index: $\ev{fx}(e, m, x)^j$ is
the $j$-th effect of some attempt of method $m$ of Step execution $e$. An
event is \emph{repeated} when an earlier event of the trace has the same
source and index. $\mathit{erase}(t)$ removes the repeated events of a
trace $t$ (\lean{eraseRepeats}). A trace $t$ is a \emph{repetition} of $u$
when $t$ is $u$ with inserted events, each equal in source and label to an
earlier event of $u$ (\lean{IsRepetition}).

\paragraph{First occurrences come in order.}
The theorem rests on one fact: for each method, the first occurrences of
the indexes $1, 2, 3, \ldots$ appear in this order, with the labels of one
effect list, even when zombies overlap with newer attempts. Three facts
give it.
\begin{enumerate}
\item Every attempt computes the same effect list and response. The
  handler is a function, and its arguments are equal across attempts: the
  snapshot is retained on every retry path (an ordinary failure, a restart
  or takeover interrupt, a returned attempt, a retry fire, the next start);
  the Step input does not change while a method is live; and H3 makes the
  attempt context irrelevant.
\item An attempt emits its effects in list order.
\item So when index $j$ first appears, some attempt has already emitted
  $1$ to $j-1$, and each of them has appeared.
\end{enumerate}
The first occurrences are therefore the effects of one \emph{virtual} task
of Dex, and every other $\ev{fx}$ event of the method equals the first
occurrence of its index. D2 is needed for the first fact: under ASYNC a
crash can lose a start that a Worker already received, and the rerun takes
a new snapshot, so its effect list can differ.

\subsection{The abstraction map}
\label{sec:abstraction}

The proof is a forward simulation of DexF by Dex under a function
$\alpha$ from DexF configurations to Dex configurations
(\lean{abstractConfig}, \file{Failure/Abstraction.lean}). Every failure
transition maps to a stutter, every other transition to Dex transitions
with the same events, and the first occurrence of each effect index to a
Worker effect of the virtual task.

$\alpha$ forgets attempt numbers, task IDs, heartbeat checkpoints, retry
timers, $Z$, $I$, and whether the server is up. It decides by the snapshot,
not by the phase: an invocation that holds a snapshot, also in Backoff or
Pending after a failure, maps to a running invocation with one virtual
task; an invocation without one never started and maps to a pending one.
This is sound because a retrying invocation looks to apply exactly like a
running one: same snapshot, same input, same locals.

The virtual task cannot be read off a live task. The attempt that went
furthest may be gone, and two runs can reach the same durable state with
different histories, for example an attempt that returned an empty effect
list and was then ignored after a failure, against one that failed before
it ran. So DexF keeps \emph{ghost} history variables (\lean{Ghost}). In
the specification no rule reads them. In Lean the step function reads them
only to check the classification Booleans that a label carries, such as
whether an effect is a first occurrence, and to update them; $\alpha$ is
their only other reader:
\begin{itemize}
\item $h(e, m)$, the largest effect index any attempt emitted, and
  $v(e, m)$, the furthest phase any attempt reached: not run, ran with a
  given effect list and response, or returned. Both are monotone. The
  virtual task has emitted the first $h(e, m)$ effects and is in phase
  $v(e, m)$.
\item $\mathit{par}(e, m)$, the handler's arguments at the first start. A
  Flow close clears an invocation and releases its snapshot while a zombie
  may still emit its first effect. $\alpha$ then keeps a canceled virtual
  task built from $\mathit{par}$, $v$, and $h$, exactly like a canceled task
  of Dex, until the last real task of the method is gone.
\item $n$ and $\zeta$: failure commands take run sequence numbers that Dex
  never sees. $\zeta$ maps the sequence number of the $k$-th applied normal
  command to $k$, and $n$ counts them, so $\alpha$ relabels every
  sequence-derived value of the model, such as the enqueue order of
  messages. $\zeta$ must be state, not a function of it: two runs can reach
  equal states, one through a normal Step start and one through a
  heartbeat at the same sequence number, and need different relabelings.
\end{itemize}

The specification classifies every command by apply's outcome. A command
takes a run sequence number exactly when it advances the run's persisted
sequence counter: normal and failure commands do, bookkeeping commands do
not. \emph{Normal} commands map to Dex commits;
\emph{failure}, \emph{bookkeeping}, internal ignored, and post-close
housekeeping commands map to stutters; a rejected client request maps to
the same Dex rejection with the same event. Worker steps follow the ghosts:
a first occurrence maps to a step of the virtual task and a repeat to a
stutter. When the last real task of a canceled invocation leaves, its drop
maps to Dex \rulename{WORKER-ABANDON} and its ignored result, if the
method had returned, to Dex \rulename{STEP-COMPLETE-STALE}; both are
progress steps. The simulation needs seven invariants (\lean{Invariant},
\file{Failure/Relation.lean}): one committed result per invocation; one
effect list per method; every task has emitted a prefix of length at most
$h$; retained state belongs to invocations that started; the abstract
state equals the Dex state up to relabeling; every Backoff invocation is
pending and has its retry timer armed, while stale retry timers may
remain; and interrupts name minted tasks.

\subsection{The theorems in Lean}
\label{sec:failure-theorems}

The trace sets are \lean{IsFailureFiniteTrace} and
\lean{IsFailureFairMaximalTrace} (DexF runs that satisfy H1 and H4, for an
admissible outcome word, and for liveness D6, HP, and the fairness of
Table~\ref{tab:failure-hypotheses}), and \lean{IsFreeFiniteTrace} for Dex
(\file{Failure/Observation.lean}). A DexF word may contain StartFlow
retries; $w{\downarrow}$ (\lean{eraseRetries}) deletes them. The statements
are in \file{Failure/Statements/Main.lean}.

\paragraph{Safety.}
Erasing the repeated effects of a finite DexF trace gives a finite Dex
trace for the word without retries:
\begin{Verbatim}[fontsize=\scriptsize]
theorem safety_failure_forward (program : Program σ)
    (hasNoFailurePolicy : HasNoFailurePolicy program) (isSyncOnly : IsSyncOnly program)
    (ignoresAttemptContext : IgnoresAttemptContext program) (environment : Dex.Environment σ)
    (requests : List (Request σ)) (trace : List (Obs σ))
    (isTrace : IsFailureFiniteTrace program environment requests trace) :
    IsFreeFiniteTrace program.failureFreeFlow environment (eraseRetries requests)
      (eraseRepeats trace)
\end{Verbatim}
\lean{safety_failure_repetition} adds that the DexF trace, with the ghost
indexes dropped, is a repetition of that Dex trace.
\lean{freeFiniteTrace_isFailureFiniteTrace} is the converse inclusion
$\mathit{Tr}_\mathit{free} \subseteq \mathit{Tr}_\mathit{fail}$: every Dex
trace is a DexF trace with no failure. This is not trivial, because DexF
carries more state. It is the only backward statement for finite traces.
The converse of the forward statement is false: not every repetition is
reachable. An attempt runs from the start and emits its effects in order,
so index $j$ never has more copies than index $j-1$. For an effect list
$[a, b, c]$, one extra $c$ without an extra $a$ and $b$ before it is a
repetition that no DexF run produces.

\paragraph{Liveness.}
A DexF step is a progress step exactly when its image under the transition
map is a Dex progress step (\lean{isProgress_eq_dexProgress}). Failures,
stale results other than that of the last task of a canceled invocation,
StartFlow retries, and repeated effects are therefore not progress, and
infinitely many of them do not make a run divergent.

\begin{Verbatim}[fontsize=\scriptsize]
theorem liveness_failure (program : Program σ)
    (hasNoFailurePolicy : HasNoFailurePolicy program) (isSyncOnly : IsSyncOnly program)
    (ignoresAttemptContext : IgnoresAttemptContext program) (environment : Dex.Environment σ)
    (word : Sequence (Request σ)) (isAdmissible : IsAdmissibleOutcomeWord environment word)
    (erasedWord : Sequence (Dex.Request σ)) (isErasure : IsRetryErasure word erasedWord)
    (trace : Sequence (Dex.Obs σ)) (marker : Marker) :
    ((∃ failureTrace, IsFailureFairMaximalTrace program word failureTrace marker ∧
        IsRepeatErasure failureTrace trace) →
      (Dex.system program.failureFreeFlow).IsFairMaximalTrace Dex.IsFair some erasedWord trace
        marker) ∧
    (HasNoRetry word →
      (Dex.system program.failureFreeFlow).IsFairMaximalTrace Dex.IsFair some erasedWord trace
        marker →
      ∃ failureTrace, IsFailureFairMaximalTrace program word failureTrace marker ∧
        IsRepeatErasure failureTrace trace)
\end{Verbatim}

The forward inclusion holds for every admissible outcome word, markers
included. The backward inclusion holds for words without retries, and it
can fail for some words with them. A deduplicated StartFlow retry is a log
entry, and its stamp strictly advances the partition clock. A word with
1,000,001 retries between two publications therefore forces more than a
second between them. A failure-free run can arm a one-second timer of one
Step and, before it is due, a zero-second timer of another Step that the
second publication starts, and fire the second timer first. With the
retries in between, the first timer is due first and takes the only
message both wait for, and no DexF run for that word has the failure-free
trace. The kernel checks the parts the model can state
(\file{Failure/Examples/RetryTiming.lean}): every StartFlow retry advances
the commit stamp, so $n$ retries put at least $n$ microseconds between two
publications, and in a failure-free Flow a delay of 1,000,001 microseconds
changes the trace. The last step, a statement about every DexF run of a
word with a million labels, remains an argument by hand. The specification
was corrected to the same two inclusions after the review of the Lean
statements.

\paragraph{End to end.}
\lean{endToEnd_safety} and \lean{endToEnd_liveness} compose these with
\lean{expressiveness_of_isWellFormed}. For every well-formed \Lang{}
program $S$, every \Lang{} trace of $S$ is the projection of a finite DexF
trace of \tr{S}, and the projection of every finite DexF trace of \tr{S} is
an \Lang{} trace of $S$ with inserted $\ev{out}(x)$ events, each equal in
value to an earlier $\ev{out}(x)$:
\begin{Verbatim}[fontsize=\scriptsize]
theorem endToEnd_safety (program : L.Program) (isWellFormed : program.IsWellFormed) :
    (∀ lTrace, IsLFiniteTrace program lTrace →
      ∃ requests failureTrace,
        IsFailureFiniteTrace (Program.ofFlow ⟦program⟧) environment requests failureTrace ∧
        failureTrace.filterMap (fun event => project event.toDex) = lTrace) ∧
    (∀ requests failureTrace,
      IsFailureFiniteTrace (Program.ofFlow ⟦program⟧) environment requests failureTrace →
      ∃ lTrace, IsLFiniteTrace program lTrace ∧
        IsOutRepetition (failureTrace.filterMap (fun event => project event.toDex)) lTrace)
\end{Verbatim}
The projection drops the source of an effect, so on the \Lang{} side an
inserted $\ev{out}(x)$ is only known to equal some earlier $\ev{out}(x)$.
For fair maximal traces the statement is stronger: the projection of
$\mathit{erase}(t)$ is a fair maximal \Lang{} trace with the same marker,
forward for every admissible outcome word and backward for words without
retries.

\paragraph{What the kernel checked.}
The proof contains no \texttt{sorry} and no \lean{native_decide}. For
\lean{safety_failure_forward}, \lean{safety_failure_repetition},
\lean{freeFiniteTrace_isFailureFiniteTrace}, \lean{liveness_failure},
\lean{endToEnd_safety}, \lean{endToEnd_liveness}, and the leaf lemmas
\lean{step_simulated}, \lean{invariant_step}, and
\lean{effectHistory_of_exec}, \texttt{\#print axioms} reports
\lean{propext}, \lean{Classical.choice}, and \lean{Quot.sound};
for \lean{isProgress_eq_dexProgress} it reports \lean{propext} and
\lean{Quot.sound}. The statements of the failure theorem rest on about
1,900 lines of new definitions that a reader must check against the
specification (\file{Failure/Config.lean}, \file{Rules.lean},
\file{Repeats.lean}, \file{Observation.lean}, and the main statements), on top
of the Dex model. The abstraction map and the invariants are internal to
the proof and need no trust.

\subsection{The proof and its work packages}
\label{sec:failure-proof}

The proof followed the method of \S\ref{sec:engineering}. The failure
model was written against the code first, and repeated adversarial
reviews checked it; its review log has 37 entries. Each finding was
checked against the code and accepted. Among them were rejected client
requests that had been made silent, a purge of closed runs that lets a
late StartFlow retry start a second run, an abstraction that was not a
function of the configuration, and a canceled invocation whose zombie
could still emit its first effect. The Lean
definitions and statements were then fixed, reviewed for vacuity and escape
hatches, and tested with an inductiveness fuzz before any proof. The proof
was cut into seven work packages with fixed statements
(Table~\ref{tab:failure-packages}), proved in parallel. No package changed a
definition or a statement.

\begin{table}[t]
\centering
\small
\begin{tabularx}{\textwidth}{@{}P{1.2cm}LP{5.2cm}@{}}
\toprule
Package & Content & Statements \\
\midrule
WP-F1 & A normal step maps to its Dex counterpart. &
  \lean{normalStep_simulated} \\
WP-F2 & A failure step is a stutter under $\alpha$. &
  \lean{failureStep_isStutter} \\
WP-F3 & First occurrences map to the virtual task, repeats to stutters;
  the ghost $h$ and the trace agree. & \lean{workerStep_simulated},
  \lean{effectHistory_of_exec} \\
WP-F4 & The invariants are preserved. & \lean{invariant_step} \\
WP-F5 & Safety, both inclusions. & \lean{safety_failure_forward},
  \lean{freeFiniteTrace_isFailureFiniteTrace}, \ldots \\
WP-F6 & The progress rule and liveness. &
  \lean{isProgress_eq_dexProgress}, \lean{liveness_failure} \\
WP-F7 & The end-to-end corollary. & \lean{endToEnd_safety},
  \lean{endToEnd_liveness} \\
\bottomrule
\end{tabularx}
\caption{The work packages of the failure proof.}
\label{tab:failure-packages}
\end{table}

\paragraph{Inductiveness fuzz.}
The fuzz (\lean{failure-fuzz}) walks pseudo-random DexF runs with failures
on a hand-written Flow and on the translations of eleven example \Lang{}
programs. At every step it checks the step simulation, the progress rule,
and a Boolean mirror of the invariants. It then perturbs every fifth
configuration in one place, keeps the copies that the mirror accepts, and
checks every enabled step from them. With the arguments
\texttt{300 5 10 12} it takes 31,160 walk steps with 194 crashes, 324
attempt losses, 189 starts on a retained snapshot, 27 repeated effects, and
118 StartFlow retries; it checks 8,425,297 single steps from 841,930 kept
copies and reports no failure. A planted program that breaks H3 makes it
report 404 failures. Three gaps in the first invariants were closed before
any proof. The fuzz found the first: an interrupt could name a task ID not
yet minted. Reviews found the other two, a Step execution ID reused when a
Step counter was out of date, and a closed run that kept a running
invocation; the extended fuzz now reproduces both.

\paragraph{Non-vacuity.}
Witnesses show that the hypotheses of \lean{safety_failure_repetition}
and of the forward inclusion of \lean{liveness_failure} hold together for
runs with crashes, retries, and zombie effects, and apply those statements
to them (\file{Failure/Examples/Nonvacuity.lean}). Each asserts
that its schedule is executable and that every step is simulated and keeps
the invariants. The backward liveness inclusion and the end-to-end
statements have no witness of their own.
\begin{itemize}
\item \lean{safetyRepetition_isNonvacuous}: a run of the translation of
  $\kw{main}(x)\colon \kw{emit}(1); \kw{exit}(0)$ with a StartFlow retry,
  a lost attempt, a zombie effect, two retry fires, a crash and its
  recovery, an interrupt, a stale zombie result, and a zombie drop. Its
  trace has three $\ev{fx}$ events, and $\mathit{erase}$ keeps one. It
  projects to the \Lang{} trace $\ev{begin}(7), \ev{out}(1),
  \ev{exit}(0)$ with two inserted $\ev{out}(1)$.
\item \lean{livenessForward_isNonvacuous_terminated}: the same run
  followed by ticks, a fair maximal DexF trace with marker \mterm{} for a
  word with a retry.
\item \lean{livenessForward_isNonvacuous_deadlocked}: the translation of
  $\kw{main}(x)\colon \kw{emit}(1); x := \kw{recv}(5)$ with a crash, an
  interrupt, a lost WaitFor attempt, two retry fires, and two stale zombie
  results, then blocked forever: a fair maximal DexF trace with marker
  \mdead{}.
\end{itemize}
A finite schedule followed by ticks forever is a fair maximal run when
three Boolean checks pass (\lean{prefixThenTicks_isFailureFairMaximalTrace}).
A further kernel-checked run takes the order the model allows and the fuzz
never reached: a zombie's result commits after a crash and before its
interrupt, which apply then ignores.

\paragraph{Counterexamples.}
Counterexamples show that H1 and H3 are needed and that a re-sent
publication counts as a second request
(\file{Failure/Examples/HypothesisCounterexamples.lean}). Without H3
(\lean{attemptSensitive_isNotRepetition}), attempt $a$ emits $a$ copies of
$a$; the retry emits $\ev{fx}(2)$, which no failure-free run emits, so
both safety statements fail.
Without H1 (\lean{exhaustion_isNotFailureFree}), an exhausted attempt
closes the Flow as Failed, which no failure-free run of that Flow does,
because it never decides ForceFail. The run of
\lean{nonRetryableFailure_isNotFailureFree} breaks H4 and also H1: in the
model a non-retryable failure is an exhaustion, so H1 implies H4, and no
counterexample shows H4 needed on its own. A re-sent
publication (\lean{resentPublication_isTwoRequests}) gives two
$\ev{pub}$ events, which no failure-free trace of the one-publication word
has, even as a repetition. H2, SYNC, and HP hold by construction in the
model, and the specification argues from the code why each is needed. All
witnesses and counterexamples depend on a subset of the three standard
axioms.

\paragraph{Size.}
Table~\ref{tab:failure-size} gives the size of the failure development. It
adds 88 files and 33,928 lines to the development of \S\ref{sec:size}, and
the specification grew to 3,534 lines. By their git author dates, the 88
non-merge commits of this work under \file{formal/} (95 with the 7 merges
of work-package branches), from the first draft of the failure model to
the last witness, fall between 10:58 and 16:16 on one day.

\begin{table}[t]
\centering
\small
\begin{tabular}{@{}lrrl@{}}
\toprule
Part of \file{Failure/} & Files & Lines & Content \\
\midrule
definitions & 6 & 2,268 & configurations, rules, repetition, $\alpha$,
  invariants, trace sets \\
\texttt{Statements/} & 2 & 391 & the leaf and main statements \\
\texttt{Proofs/} & 49 & 18,940 & WP-F1 to WP-F4 and the leaf part of
  WP-F6 \\
\texttt{Composite/} & 17 & 8,305 & WP-F5 to WP-F7 \\
\texttt{Examples/} & 9 & 3,449 & kernel-checked runs, witnesses,
  counterexamples, review probes, the fuzz \\
tooling & 5 & 575 & decidable equality, checkers, aggregators \\
\midrule
Total & 88 & 33,928 & \\
\bottomrule
\end{tabular}
\caption{The failure development at commit \texttt{dd449bd57}, counted
with \texttt{wc -l}. \file{FailureFuzz.lean}, the fuzz entry point, adds
59 lines outside \file{Failure/}.}
\label{tab:failure-size}
\end{table}
